\documentclass[11pt,english]{article}
\usepackage[T1]{fontenc}
\usepackage[utf8]{inputenc}
\usepackage{geometry}
\geometry{verbose,tmargin=1in,bmargin=1in,lmargin=1in,rmargin=1in}
\usepackage{amsmath}
\usepackage{amssymb}
\usepackage{bm}
\usepackage{booktabs}
\usepackage{enumitem}
\setlength{\parskip}{\smallskipamount}
\setlength{\parindent}{0pt}

\input{Macros.tex}

\title{Problem Set 1 --- Solutions\\
\large Labor Economics (Saggio), Summer 2026}
\author{Pablo Fiorentini}
\date{}

\begin{document}
\maketitle

%=======================================================================
\section*{Question 1: Returns to Schooling}
%=======================================================================

An infinitely-lived agent picks schooling length $S$, faces discount rate $r$,
earns the scholarship flow $T$ on $[0,S]$, and earns $y(S)$ forever after date
$S$. She maximises the present value of earnings plus transfers.

\subsection*{Part 1: Objective and optimality condition}

In continuous time the objective is
\begin{equation}
V(S)=\underbrace{\int_{0}^{S}T\,e^{-rt}\,dt}_{\text{scholarship while in school}}
     +\underbrace{\int_{S}^{\infty}y(S)\,e^{-rt}\,dt}_{\text{earnings after school}}
     =\frac{T\left(1-e^{-rS}\right)}{r}+\frac{y(S)\,e^{-rS}}{r}.
\end{equation}
Differentiating with respect to $S$,
\begin{equation}
V'(S)=T e^{-rS}+\frac{y'(S)e^{-rS}-r\,y(S)e^{-rS}}{r}
     =e^{-rS}\!\left[T+\frac{y'(S)}{r}-y(S)\right].
\end{equation}
Setting $V'(S)=0$ and cancelling $e^{-rS}>0$ gives the optimality condition
\begin{equation}\boxed{\;y'(S^{*})=r\bigl(y(S^{*})-T\bigr).\;}\end{equation}

\textbf{Interpretation.} Rewrite as $\dfrac{y'(S^{*})}{r}=y(S^{*})-T$. The
left-hand side is the \emph{marginal benefit} of one more year of school: the
permanent earnings increment $y'(S)$ capitalised at rate $r$. The right-hand
side is the \emph{marginal cost}: the earnings $y(S)$ forgone by staying in
school one more instant, \emph{net} of the scholarship $T$ that is still
collected. The agent keeps studying until the discounted value of the higher
earnings path just balances the net opportunity cost of delaying entry into the
labor market. A more generous scholarship $T$ lowers the marginal cost and
raises optimal schooling.

\subsection*{Part 2: Ability and the optimal duration of schooling}

With $y_{i}(S)=\exp(\alpha_{i}+\beta S)$ we have $y_{i}'(S)=\beta\,y_{i}(S)$.
The optimality condition becomes
\begin{equation}
\beta\,y_{i}=r\bigl(y_{i}-T\bigr)
\quad\Longrightarrow\quad
y_{i}(S_{i}^{*})=\frac{rT}{r-\beta},
\end{equation}
which requires $r>\beta$ (so that the cost of delay eventually dominates and an
interior optimum exists). Taking logs of
$\exp(\alpha_{i}+\beta S_{i}^{*})=rT/(r-\beta)$,
\begin{equation}\boxed{\;
S_{i}^{*}=\frac{1}{\beta}\left[\log\!\Bigl(\frac{rT}{r-\beta}\Bigr)-\alpha_{i}\right].
\;}\end{equation}

\textbf{Do higher-ability agents get more or less schooling?} \emph{Less}: the
coefficient on $\alpha_{i}$ is $-1/\beta<0$. Intuitively, the marginal return to
schooling $\beta$ is common to everyone, but the optimal stopping rule pins the
\emph{level} of earnings at $y^{*}=rT/(r-\beta)$, identical across agents. A
high-ability agent starts from a higher $\alpha_{i}$, so she reaches that target
earnings level with fewer years of school. Equivalently, ability and schooling
are substitutes in producing $y$, and the opportunity cost of an extra year
(forgone earnings $y_i(S)-T$) rises faster for the high-ability agent, pushing
her out of school sooner.

\subsection*{Part 3: OLS bias with heterogeneous ability and scholarships}

Now $\alpha_{i}$ and $T_{i}$ vary across agents while $\beta$ is common. From
Part 2, at the optimum $y_{i}=rT_{i}/(r-\beta)$, so
\begin{align}
\log y_{i} &= \log\!\Bigl(\tfrac{r}{r-\beta}\Bigr)+\log T_{i}, \\[2pt]
S_{i}      &= \frac{1}{\beta}\left[\log\!\Bigl(\tfrac{r}{r-\beta}\Bigr)+\log T_{i}-\alpha_{i}\right].
\end{align}
Note $\log y_i$ depends on $\log T_i$ only (not on $\alpha_i$ directly): ability
heterogeneity is ``undone'' in equilibrium earnings because high-$\alpha$ agents
choose less schooling. The OLS slope of $\log y$ on $S$ has probability limit
$\plim\hat{b}=\Cov(\log y_{i},S_{i})/\Var(S_{i})$. Using
$\sigma_{T}^{2}=\Var(\log T_i)$, $\sigma_{\alpha}^{2}=\Var(\alpha_i)$,
$\sigma_{T\alpha}=\Cov(\log T_i,\alpha_i)$:
\begin{align}
\Cov(\log y_i,S_i)&=\Cov\!\Bigl(\log T_i,\tfrac{1}{\beta}(\log T_i-\alpha_i)\Bigr)
=\tfrac{1}{\beta}\bigl(\sigma_T^2-\sigma_{T\alpha}\bigr),\\
\Var(S_i)&=\tfrac{1}{\beta^2}\Var(\log T_i-\alpha_i)
=\tfrac{1}{\beta^2}\bigl(\sigma_T^2+\sigma_\alpha^2-2\sigma_{T\alpha}\bigr).
\end{align}
Therefore
\begin{equation}\boxed{\;
\plim\hat{b}=\beta\;\frac{\sigma_T^2-\sigma_{T\alpha}}
{\sigma_T^2+\sigma_\alpha^2-2\sigma_{T\alpha}}.
\;}\end{equation}

\textbf{Discussion of the bias.}
\begin{itemize}[nosep]
\item \emph{No ability heterogeneity} ($\sigma_\alpha^2=0,\ \sigma_{T\alpha}=0$):
$\plim\hat b=\beta$. OLS is consistent --- all schooling variation is driven by
the ``cost shifter'' $T_i$, which is independent of the outcome equation, so
$T_i$ acts like an instrument and there is no bias.
\item \emph{Ability variance} ($\sigma_\alpha^2>0$, $\sigma_{T\alpha}=0$):
$\plim\hat b=\beta\,\sigma_T^2/(\sigma_T^2+\sigma_\alpha^2)<\beta$. Variation in
$\alpha_i$ adds noise to $S_i$ that is unrelated to $\log y_i$, \emph{attenuating}
the slope toward zero.
\item \emph{Sorting} ($\sigma_{T\alpha}\neq0$): the covariance term shifts the
bias. Because here $\log y_i$ rises with $T_i$ but high-ability agents take
\emph{less} schooling, a positive $\sigma_{T\alpha}$ reduces the numerator and
denominator in offsetting ways; the net effect depends on the relative
magnitudes. The key message is that OLS recovers $\beta$ only when schooling
choices are uncorrelated with ability.
\end{itemize}

\subsection*{Part 4: OLS bias with heterogeneous returns}

Now $T_{i}=T$ is constant but returns are heterogeneous,
$y_{i}(S)=\exp(\alpha_{i}+\beta_{i}S)$. The Part-1 condition
$\beta_i y_i=r(y_i-T)$ gives $y_i(S_i^*)=rT/(r-\beta_i)$, so
\begin{align}
\log y_{i} &= \log(rT)-\log(r-\beta_{i}), \label{eq:p4y}\\
S_{i}      &= \frac{1}{\beta_{i}}\bigl[\log(rT)-\log(r-\beta_{i})-\alpha_{i}\bigr].
\label{eq:p4S}
\end{align}
This is a \emph{correlated random coefficients} model: the slope $\beta_i$ that
generates earnings is the same object that enters the schooling choice. Because
\eqref{eq:p4S} is nonlinear in $\beta_i$, we obtain a closed form by a
first-order expansion around the means $\beta\equiv\E[\beta_i]$,
$\alpha\equiv\E[\alpha_i]$. Write $\beta_i=\beta+b_i$, $\alpha_i=\alpha+a_i$,
with $\Var(b_i)=\sigma_\beta^2$, $\Var(a_i)=\sigma_\alpha^2$,
$\Cov(b_i,a_i)=\sigma_{\beta\alpha}$, and define
\[
\varphi\equiv\frac{1}{r-\beta},\qquad
\bar S\equiv\frac{\log(rT)-\log(r-\beta)-\alpha}{\beta},\qquad
\psi\equiv\varphi-\bar S .
\]
Since $\partial[-\log(r-\beta_i)]/\partial\beta_i=\varphi$ at $\beta_i=\beta$,
equation \eqref{eq:p4y} linearises to $\log y_i\approx\text{const}+\varphi\,b_i$,
and \eqref{eq:p4S} to
\[
S_i-\bar S\approx\frac{1}{\beta}\bigl[(\varphi-\bar S)\,b_i-a_i\bigr]
            =\frac{1}{\beta}\bigl[\psi\,b_i-a_i\bigr].
\]
Hence
\begin{align}
\Cov(\log y_i,S_i)&\approx\frac{\varphi}{\beta}\bigl(\psi\,\sigma_\beta^2-\sigma_{\beta\alpha}\bigr),\\
\Var(S_i)&\approx\frac{1}{\beta^2}\bigl(\psi^2\sigma_\beta^2+\sigma_\alpha^2-2\psi\,\sigma_{\beta\alpha}\bigr),
\end{align}
and the OLS slope satisfies
\begin{equation}\boxed{\;
\plim\hat{b}\approx\varphi\,\beta\;
\frac{\psi\,\sigma_\beta^2-\sigma_{\beta\alpha}}
{\psi^2\sigma_\beta^2+\sigma_\alpha^2-2\psi\,\sigma_{\beta\alpha}},
\qquad \varphi=\tfrac{1}{r-\beta},\ \ \psi=\varphi-\bar S.
\;}\end{equation}
(The constants $\varphi$ and $\psi$ fix the linearisation point; the
\emph{bias} is governed by the second moments
$\sigma_\beta^2,\sigma_\alpha^2,\sigma_{\beta\alpha}$.)

\textbf{Discussion of the bias.}
\begin{itemize}[nosep]
\item \emph{Homogeneous returns} ($\sigma_\beta^2=0$, hence
$\sigma_{\beta\alpha}=0$): $\plim\hat b\to0$. With $T$ constant and a common
$\beta$, equation \eqref{eq:p4y} makes $\log y_i$ \emph{identical} across agents
($y$ is pinned to $rT/(r-\beta)$), so there is no earnings variation to explain
and the slope collapses to zero. \emph{All} earnings dispersion in this model
comes from heterogeneous returns $\beta_i$.
\item \emph{Return heterogeneity alone} ($\sigma_\beta^2>0$,
$\sigma_\alpha^2=\sigma_{\beta\alpha}=0$): $\plim\hat b\approx\varphi\beta/\psi
\neq\beta$. Even with no ability heterogeneity, OLS does \emph{not} recover the
average return $\beta$, because agents with different $\beta_i$ optimally choose
different $S_i$ --- selection on gains contaminates the slope.
\item \emph{Sorting on gains} ($\sigma_{\beta\alpha}\neq0$) and ability variance
($\sigma_\alpha^2>0$) move the estimand further from $\beta$ and, as the sign of
$\psi$ can be negative, $\plim\hat b$ may even take the \emph{opposite sign} of
the average return. The practical lesson is the same as Part 3: a bivariate OLS
of log earnings on schooling generally fails to identify the causal return when
schooling is chosen, and the failure is more severe when the return itself is
heterogeneous and correlated with the choice.
\end{itemize}

%=======================================================================
\section*{Question 2: Keane and Wolpin (1997)}
%=======================================================================

Ages $a\in\{1,2\}$; choices $m\in\{1,2,3\}=\{$white collar, blue collar,
school$\}$. Flow utilities
\[
R_{m}(a)=\alpha_{m}+\beta_{m}G(a)+\gamma_{m}X_{m}(a)+\epsilon_{m}(a)\quad(m=1,2),
\qquad R_{3}(a)=\epsilon_{3}(a),
\]
with the school alternative normalised
($\alpha_{3}=\beta_{3}=\gamma_{3}=0$). The $\epsilon_{m}(a)$ are i.i.d. Type-I
Extreme Value across alternatives and time; the period-2 shocks are unknown in
period 1. The agent enters period 1 with $G(1)=X_1(1)=X_2(1)=0$ and discounts
the future at $\delta\in(0,1)$.

\subsection*{Part 1: State vector and period-2 choice probabilities}

The state entering age 2 (before $\epsilon_m(2)$ is realised) is
\[
S(2)=\bigl(G(2),\,X_1(2),\,X_2(2)\bigr).
\]
Because the agent begins period 1 empty and makes a single choice, exactly one
state component equals $1$ at age 2. Thus $S(2)$ takes only \emph{three} values:
\[
\underbrace{(0,1,0)}_{\text{worked WC in 1}},\qquad
\underbrace{(0,0,1)}_{\text{worked BC in 1}},\qquad
\underbrace{(1,0,0)}_{\text{schooled in 1}}.
\]
Period 2 is the last period, so it is a static multinomial-logit choice. Writing
the deterministic parts
\[
v_{1}(2)=\alpha_1+\beta_1 G(2)+\gamma_1 X_1(2),\quad
v_{2}(2)=\alpha_2+\beta_2 G(2)+\gamma_2 X_2(2),\quad
v_{3}(2)=0,
\]
the Type-I EV assumption gives, for each $S(2)$,
\begin{equation}
\Pr\bigl(C(2)=m\mid S(2)\bigr)=\frac{\exp\!\bigl(v_{m}(2)\bigr)}
{\exp\!\bigl(v_{1}(2)\bigr)+\exp\!\bigl(v_{2}(2)\bigr)+\exp\!\bigl(v_{3}(2)\bigr)},
\qquad m=1,2,3 .
\end{equation}

\subsection*{Part 2: Expected value function at age 2}

With $V_{2}(S(2))=\max_{m}R_{m}(2;S(2))$ and the Type-I EV hint
$\E[\max_k U_{ik}+\eta_{ik}]=\nu+\log\sum_k\exp(U_{ik})$, the expectation
(prior to the realisation of $\epsilon_m(2)$, dropping the constant $\nu$) is
\begin{equation}
\E\bigl[V_{2}(S(2))\bigr]
=\log\!\Bigl(\exp\!\bigl(v_{1}(2)\bigr)+\exp\!\bigl(v_{2}(2)\bigr)+1\Bigr),
\end{equation}
evaluated at each of the three states (note $\exp(v_3(2))=e^0=1$). Explicitly,
\begin{align*}
\text{WC experience }(0,1,0):&\quad \log\!\bigl(e^{\alpha_1+\gamma_1}+e^{\alpha_2}+1\bigr),\\
\text{BC experience }(0,0,1):&\quad \log\!\bigl(e^{\alpha_1}+e^{\alpha_2+\gamma_2}+1\bigr),\\
\text{Education }(1,0,0):&\quad \log\!\bigl(e^{\alpha_1+\beta_1}+e^{\alpha_2+\beta_2}+1\bigr).
\end{align*}

\subsection*{Part 3: Period-1 expected payoffs and choice probabilities}

In period 1 the deterministic flows are the intercepts only (state is empty):
$v_1(1)=\alpha_1$, $v_2(1)=\alpha_2$, $v_3(1)=0$. Each choice deterministically
sets next period's state, so the choice-specific value (excluding $\epsilon_m(1)$)
is the current flow plus the discounted expected continuation value of the state
it leads to:
\begin{align}
w_{1}&=\alpha_{1}+\delta\,\E\bigl[V_2\mid(0,1,0)\bigr] && \text{(white collar)},\\
w_{2}&=\alpha_{2}+\delta\,\E\bigl[V_2\mid(0,0,1)\bigr] && \text{(blue collar)},\\
w_{3}&=0+\delta\,\E\bigl[V_2\mid(1,0,0)\bigr] && \text{(school)}.
\end{align}
Adding the i.i.d. Type-I EV shocks $\epsilon_m(1)$, the \emph{ex ante} period-1
choice probabilities are again multinomial logit:
\begin{equation}
\Pr\bigl(C(1)=m\bigr)=\frac{\exp(w_{m})}{\exp(w_{1})+\exp(w_{2})+\exp(w_{3})},
\qquad m=1,2,3 .
\end{equation}
School has no current payoff but yields the largest education continuation value
$\E[V_2\mid(1,0,0)]$; the agent trades current earnings against the discounted
option value of being more educated (or more experienced) at age 2.

\subsection*{Part 4: Likelihood}

Individual $i$ is observed choosing the sequence $(C_i(1),C_i(2))$. The
period-1 choice determines $S_i(2)$, so by the law of total probability the
individual likelihood factors as
\begin{equation}
L_i(\theta)=\Pr\bigl(C(1)=C_i(1)\bigr)\cdot
\Pr\bigl(C(2)=C_i(2)\mid S_i(2)\bigr),
\end{equation}
where $S_i(2)$ is the state implied by $C_i(1)$. Taking logs and summing over the
sample, the log-likelihood is
\begin{equation}\boxed{\;
\ell(\theta)=\sum_{i=1}^{N}\Bigl[\log\Pr\bigl(C(1)=C_i(1)\bigr)
+\log\Pr\bigl(C(2)=C_i(2)\mid S_i(2)\bigr)\Bigr],
\;}\end{equation}
with $\theta=(\alpha_1,\beta_1,\gamma_1,\alpha_2,\beta_2,\gamma_2,\delta)'$ and
the probabilities given in Parts 1 and 3.

\subsection*{Part 5: Maximum-likelihood estimation}

I estimate $\theta$ by maximum likelihood on the $N=5{,}000$ simulated choice
sequences in \texttt{kw\_example.csv}, minimising $-\ell(\theta)$ in Python
(\texttt{scipy.optimize.minimize}, BFGS --- the analogue of Matlab's
\texttt{fminunc}). I search over $\log\delta$ to keep $\delta>0$, and use several
random starting values, retaining the global optimum. The code is in
\texttt{kw\_estimate.py}; standard errors are described below.

\begin{table}[h]
\centering
\caption{Maximum-likelihood estimates of the simplified Keane--Wolpin model}
\begin{tabular}{lcccc}
\toprule
 & & \multicolumn{3}{c}{Standard error}\\ \cmidrule(l){3-5}
Parameter & Estimate & Hessian & OPG & Sandwich \\ \midrule
$\alpha_1$ (white-collar intercept)        & $0.4234$  & $0.0585$ & $0.0624$ & $0.0552$ \\
$\beta_1$ (return to education, WC)         & $0.3681$  & $0.0800$ & $0.0873$ & $0.0738$ \\
$\gamma_1$ (return to WC experience)        & $0.0607$  & $0.0844$ & $0.0881$ & $0.0814$ \\
$\alpha_2$ (blue-collar intercept)          & $0.4237$  & $0.0422$ & $0.0431$ & $0.0416$ \\
$\beta_2$ (return to education, BC)          & $0.1215$  & $0.0635$ & $0.0636$ & $0.0634$ \\
$\gamma_2$ (return to BC experience)        & $0.3513$  & $0.0792$ & $0.0823$ & $0.0768$ \\
$\delta$ (discount factor)                  & $0.8395$  & $0.3818$ & $0.4147$ & $0.3541$ \\ \midrule
\multicolumn{5}{l}{Log-likelihood $=-10{,}734.24$,\quad $N=5{,}000$}\\
\bottomrule
\end{tabular}
\end{table}

\paragraph{Standard errors (extra credit).} The MLE
$\hat\theta\overset{a}{\sim}\mathcal{N}\!\bigl(\theta_0,\,V/N\bigr)$, and I report
three consistent estimators of the asymptotic variance, all computed on the
$\rho=\log\delta$ parameterisation and then mapped to $\delta$ by the delta
method (see below). Let $s_i(\theta)=\partial\ell_i/\partial\theta$ be the
score of observation $i$ and $\hat H=-\sum_i\partial^2\ell_i/\partial\theta\,\partial\theta'$
the observed (negative) Hessian, both evaluated at $\hat\theta$:
\begin{equation}
\underbrace{\widehat V_{H}=\hat H^{-1}}_{\text{observed information}}
,\qquad
\underbrace{\widehat V_{J}=\Bigl(\textstyle\sum_i s_i s_i'\Bigr)^{-1}}_{\text{OPG / BHHH}}
,\qquad
\underbrace{\widehat V_{S}=\hat H^{-1}\Bigl(\textstyle\sum_i s_i s_i'\Bigr)\hat H^{-1}}_{\text{sandwich / robust}} .
\end{equation}
I compute $\hat H$ by a finite-difference Hessian of $-\ell$ and the scores
$s_i$ by central differences; because each observation's likelihood depends only
on its $(C(1),C(2))$ cell, the score takes only nine distinct values, so the
OPG matrix is $\sum_{\text{cells}} n_{\text{cell}}\,s\,s'$. The reported baseline
SEs are the Hessian-based $\sqrt{\operatorname{diag}\widehat V_{H}}$.

\emph{Delta method for $\delta$.} I optimise over $\rho=\log\delta$; since
$\delta=e^{\rho}$ has $\partial\delta/\partial\rho=\delta$, the SE of $\delta$ is
$\widehat{\SE}(\delta)=\delta\cdot\widehat{\SE}(\rho)$ (with
$\widehat{\SE}(\log\delta)=0.4548$ under the Hessian estimator).

\emph{Information-matrix equality.} Under correct specification the
information-matrix equality $H=J$ holds, so the three estimators should
coincide asymptotically. They do: across the six slope/intercept parameters the
Hessian, OPG, and sandwich SEs agree to within about ten percent of one another
(differing only in the second decimal), which is reassuring evidence
that the model and likelihood are correctly specified. The lone caveat is
$\delta$, whose SE ($\approx0.38$--$0.41$) is large under every method because
the likelihood is flat in $\delta$.

\textbf{Goodness of fit.} With seven parameters for the eight free cell
probabilities of the $3\times3$ table of $(C(1),C(2))$ sequences, the model
reproduces every empirical cell frequency to within $0.0014$ in probability,
confirming the likelihood is correctly specified and the optimum is the global
one.

\textbf{Economic interpretation.}
\begin{itemize}[nosep]
\item \emph{Intercepts.} $\alpha_1\approx\alpha_2\approx0.42$: for an
inexperienced, uneducated entrant, the two occupations offer almost identical
baseline utility, so the period-1 choice is driven largely by their differing
\emph{continuation} values.
\item \emph{Returns to education.} $\beta_1=0.37\gg\beta_2=0.12$: a year of
schooling raises white-collar utility about three times as much as blue-collar
utility. Education is the human-capital investment that pays off mainly in the
white-collar track ($\beta_1$ is precisely estimated; $\beta_2$ is only
marginally distinguishable from zero).
\item \emph{Returns to experience.} The pattern reverses for on-the-job
experience: $\gamma_2=0.35$ (precise) versus $\gamma_1=0.06$ (indistinguishable
from zero). Blue-collar work builds valuable occupation-specific human capital,
whereas white-collar experience adds little.
\item \emph{Discount factor.} $\hat\delta=0.84\in(0,1)$ is economically
reasonable, but its standard error is large ($0.38$): the log-likelihood is flat
in $\delta$ over roughly $[0.7,1.0]$, so $\delta$ is only weakly identified from
two periods of choices. (The BFGS routine's convergence flag reflects this flat
ridge; the gradient norm at the reported optimum is $\sim\!2\times10^{-4}$.)
\end{itemize}

\textbf{The period-1 trade-off.} Schooling delivers \emph{no} current payoff
($R_3=\epsilon_3$) but is the only way to acquire education $G(2)=1$, which
sharply raises age-2 white-collar utility (via $\beta_1=0.37$) and modestly
raises blue-collar utility ($\beta_2=0.12$). Blue-collar work pays the immediate
intercept $\alpha_2$ \emph{and} builds high-return experience ($\gamma_2=0.35$),
making it attractive on both margins. White-collar work pays $\alpha_1$ now but
accumulates little useful experience ($\gamma_1\approx0$), so its appeal in
period 1 rests mostly on the current flow. The agent therefore weighs immediate
utility against the discounted option value of human capital, and how strongly
that option value pulls her toward school is governed by $\delta$.

\end{document}
